\chapter{Doğrulama, Güvenilirlik ve Operasyonel Kullanım}
\label{ch:verification_reliability_operations}

Önceki bölümlerde yerçekimi modellemesinin teorik altyapısı, Ay çevresi dinamiğinin sayısal propagasyonu, yüksek sadakatli kuvvet modelleri ve ST-LRPS artık potansiyel vekil modeli ayrıntılı biçimde ele alınmıştır. Ancak yüksek sadakatli bir simülasyon altyapısında yalnızca modelin kurulmuş olması yeterli değildir. Modelin hangi veriyle çalıştığı, hangi varsayımları yaptığı, hangi koşullarda güvenilir olduğu ve hangi sonuçların operasyonel karar için yeterli kabul edileceği de aynı açıklıkta tanımlanmalıdır.

Bu bölüm, çalışmanın eksik kalan mühendislik güvenilirliği tarafını tamamlamak için eklenmiştir. Amaç; \texttt{LUNAR\_SIMULATION} projesinin yalnızca çalışan bir yazılım değil, izlenebilir, tekrar üretilebilir, test edilebilir ve görev kararlarını destekleyebilir bir analiz platformu olarak nasıl ele alınması gerektiğini göstermektir.

\section{Doğrulama ve Geçerleme Ayrımı}
\label{sec:vv_difference}

Uzay görevi simülasyonlarında sık yapılan hata, ``kod çalışıyor'' ifadesini ``model doğrudur'' ifadesiyle karıştırmaktır. Bu iki ifade farklı doğrulama katmanlarına karşılık gelir.

\begin{itemize}
    \item \textbf{Doğrulama (verification):} Yazılımın, yazılması hedeflenen matematiksel modeli doğru biçimde çözüp çözmediğini inceler. Örneğin iki-cisim probleminde enerji korunumu, analitik \(J_2\) precession testi veya zaman dönüşümü round-trip kontrolü bu sınıfa girer.
    \item \textbf{Geçerleme (validation):} Matematiksel modelin fiziksel gerçekliği yeterince temsil edip etmediğini inceler. Örneğin GRAIL tabanlı bir Ay yerçekimi modeliyle propagasyon sonucunun bilinen mascon etkilerini üretmesi veya LOLA topografyasına dayalı çarpışma analizinin yüzey morfolojisiyle tutarlı olması bu sınıfa girer.
\end{itemize}

Bu nedenle proje için doğrulama stratejisi katmanlı kurulmalıdır:
\[
\begin{aligned}
\text{birim test}
&\longrightarrow
\text{model testi}
\longrightarrow
\text{entegrasyon testi} \\
&\longrightarrow
\text{senaryo testi}
\longrightarrow
\text{operasyonel kabul}.
\end{aligned}
\]

Bu zincirde bir üst katmanın güvenilirliği, alt katmanların sessiz hata üretmemesine bağlıdır. Örneğin ST-LRPS modeli iyi bir validasyon RMSE değeri verse bile, veri kümesinin yanlış merkezi cisim metadata'sı taşıması durumunda sonuç fiziksel olarak geçersizdir. Bu yüzden sistemde fail-fast metadata kontrolleri, yalnızca yazılım kalitesi değil doğrudan fizik güvenilirliği gereğidir.

\section{Katmanlı Test Stratejisi}
\label{sec:test_strategy}

Projedeki test stratejisi tek bir büyük uçtan uca test yerine, farklı hata sınıflarını yakalayan küçük ama anlamlı test kümelerine dayanmalıdır. Böylece hata kaynağı daha hızlı izole edilir.

\subsection{Birim Testleri}

Birim testleri, modüllerin en küçük hesaplama sözleşmelerini denetler. Bu testler hızlı çalışmalı ve mümkün olduğunca deterministik olmalıdır. Önerilen birim test başlıkları şunlardır:

\begin{itemize}
    \item \textbf{Zaman dönüşümleri:} UTC, J2000 saniyesi ve ISO format dönüşümleri mikro-saniye sınırlarında round-trip test edilmelidir.
    \item \textbf{Durum dönüşümleri:} Kartezyen durum vektörü ile klasik orbital elemanlar arasında dönüşüm yapılıp geri dönüldüğünde hata tolerans içinde kalmalıdır.
    \item \textbf{Yerçekimi katsayı okuma:} GFC/SH dosyasından okunan \(\mu\), referans yarıçapı ve derece bilgisi metadata ile tutarlı olmalıdır.
    \item \textbf{Topografya örnekleme:} Sabit küre sağlayıcısı ve LDEM sağlayıcısı aynı arayüzü sunmalı; koordinat sınırlarında örnekleme NaN üretmemelidir.
    \item \textbf{ST-LRPS scaler:} Konum ölçeği Ay merkezini kaydırmamalı, yani \(\mathbf{x}_0=\mathbf{0}\) sözleşmesi korunmalıdır.
\end{itemize}

\subsection{Model Testleri}

Model testleri, fiziksel alt modellerin beklenen davranışlarını ayrı ayrı inceler. Örneğin yalnızca merkezi Ay alanı açıkken iki-cisim enerjisi yaklaşık korunmalıdır. Sferik harmonik derece artırıldığında ivme farkı büyüklük olarak makul ve sürekli değişmelidir. Üçüncü cisim etkisi açıldığında ivme, mutlak Güneş veya Dünya çekiminden ziyade diferansiyel perturbasyon formunda hesaplanmalıdır.

Bu test sınıfında temel kabul, modelin belirli sınır koşullarında analitik veya yüksek güvenilirlikli referansla uyumlu olmasıdır. Örneğin:
\[
\epsilon_E(t)
=
\frac{|E(t)-E(0)|}{|E(0)|}
\]
iki-cisim testlerinde uzun süre boyunca küçük kalmalıdır. Buradaki hedef değer integratör toleransına, adım sınırına ve süreye bağlıdır; fakat hata trendi monoton büyüyen bir sayısal drift gösteriyorsa yalnızca anlık tolerans değeri yeterli kabul edilmemelidir.

\subsection{Entegrasyon Testleri}

Entegrasyon testleri, modüllerin birlikte çalışmasını denetler. Bu sınıfta amaç tek bir fonksiyonun doğru olup olmadığını değil, veri akışının uçtan uca tutarlı olup olmadığını görmektir. Örneğin:

\begin{itemize}
    \item UI tarafında seçilen yerçekimi backend değerinin \texttt{main.py} ve \texttt{mc\_runner.py} tarafına aynı anlamla taşınması,
    \item \texttt{--surrogate-gravity-model-dir} ile verilen ST-LRPS klasörünün config, scaler ve checkpoint dosyalarıyla birlikte yüklenebilmesi,
    \item topografya klasörü boşsa sistemin sabit yarıçap sağlayıcısına geçtiğini loglarda açıkça göstermesi,
    \item Monte Carlo çalışması bittiğinde terminal ve süreç durumunun yeniden koşu başlatmaya izin verecek şekilde kapanması
\end{itemize}
entegrasyon testlerinin kapsamına girer.

\subsection{Senaryo Testleri}

Senaryo testleri, gerçek kullanıcı akışlarını temsil eder. Bir görev konfigürasyonu seçilir, simülasyon çalıştırılır, çıktı klasörü üretilir, rapor oluşturulur ve gerekiyorsa Monte Carlo analizi yapılır. Bu testler daha yavaş olabilir; fakat yazılımın operasyonel kullanıma hazır olup olmadığını en iyi bu testler gösterir.

Önerilen senaryo seti:
\begin{itemize}
    \item \textbf{Kısa klasik SH koşusu:} 1 günlük, düşük dereceli klasik sferik harmonik test.
    \item \textbf{Uzun klasik SH koşusu:} 30--100 günlük, tolerans ve adım politikası testi.
    \item \textbf{Topografya çarpışma koşusu:} Alçak irtifalı, LDEM verisiyle yüzey temas testi.
    \item \textbf{ST-LRPS karşılaştırma koşusu:} Aynı başlangıç koşulunda klasik SH ve ST-LRPS fark analizi.
    \item \textbf{Monte Carlo koşusu:} Küçük senaryo sayısıyla hızlı duman testi, büyük senaryo sayısıyla istatistiksel kararlılık testi.
\end{itemize}

\section{Hata Bütçesi ve Güvenilirlik Kaynakları}
\label{sec:error_budget}

Ay yörüngesi propagasyonunda toplam hata tek bir kaynaktan gelmez. Sayısal integratör hatası, yerçekimi modeli truncation hatası, efemeris interpolasyon hatası, referans çerçevesi dönüşümü, topografya çözünürlüğü, yüzey modeli varsayımı, uzay aracı parametre belirsizliği ve ST-LRPS model hatası birlikte değerlendirilmelidir.

Toplam hata kavramsal olarak şu şekilde ayrıştırılabilir:
\[
\epsilon_{\mathrm{total}}
\approx
\epsilon_{\mathrm{num}}
+ \epsilon_{\mathrm{grav}}
+ \epsilon_{\mathrm{eph}}
+ \epsilon_{\mathrm{frame}}
+ \epsilon_{\mathrm{surface}}
+ \epsilon_{\mathrm{sc}}
+ \epsilon_{\mathrm{surrogate}}.
\]
Bu ifade kesin bir eşitlik değildir; hata kaynakları bağımsız olmayabilir. Ancak hangi bileşenin görev sonucunu ne kadar etkilediğini düşünmek için yararlı bir mühendislik çerçevesi sunar.

\begingroup
\footnotesize
\begin{longtable}{>{\raggedright\arraybackslash}p{0.22\textwidth}
                  >{\raggedright\arraybackslash}p{0.32\textwidth}
                  >{\raggedright\arraybackslash}p{0.36\textwidth}}
\toprule
\textbf{Hata Kaynağı} & \textbf{Tipik Belirti} & \textbf{Kontrol / Azaltma Yöntemi} \\
\midrule
\endhead
Sayısal integrasyon &
Enerji drift'i, adım küçülmesi, tolerans hassasiyeti. &
DOP853 tolerans taraması, maksimum adım sınırı, iki-cisim ve düşük dereceli referans testleri. \\
Yerçekimi truncation &
Alçak irtifada faz ve irtifa zarfı farkı. &
Derece taraması, SH20/SH50/SH100 karşılaştırması, ST-LRPS artık alan testi. \\
Efemeris / zaman &
Faz kayması, üçüncü cisim ivmesinde tutarsızlık. &
UTC-J2000 round-trip, efemeris tablo aralığı kontrolü, interpolasyon sınır testi. \\
Referans çerçevesi &
Ay-sabit alanın yanlış eksende uygulanması. &
Kuaterniyon norm kontrolü, ters dönüşüm testi, fixed/inertial dönüşüm round-trip. \\
Topografya &
Yanlış çarpışma zamanı veya negatif irtifa. &
LDEM label/IMG eşleşme kontrolü, sabit küre karşılaştırması, lokal yüzey yarıçapı testi. \\
Uzay aracı parametreleri &
SRP/albedo hassasiyetinde büyük saçılım. &
Alan/kütle, \(C_D\), \(C_R\) belirsizliklerini Monte Carlo ile örnekleme. \\
ST-LRPS model hatası &
Belirli irtifa kabuğunda ivme yön hatası veya OOD bozulması. &
Bağımsız test, OOD tablosu, radial/cross-radial hata ve shell-wise değerlendirme. \\
\bottomrule
\end{longtable}
\endgroup

Hata bütçesi özellikle rapor yorumlarında önemlidir. Örneğin ST-LRPS ve klasik SH sonuçları arasında küçük bir fark görülmesi, otomatik olarak ST-LRPS'in hatalı olduğu anlamına gelmez; klasik modelin seçilen derece truncation hatası veya farklı maksimum adım politikası da bu farkı üretebilir. Bu yüzden karşılaştırmalı koşularda aynı başlangıç durumu, aynı zaman aralığı, aynı çıktı örnekleme aralığı ve aynı tolerans politikası kullanılmalıdır.

\section{Veri İzlenebilirliği ve Tekrar Üretilebilirlik}
\label{sec:data_provenance}

Projenin güvenilirliği, yalnızca hesaplama fonksiyonlarının doğruluğuna değil, kullanılan veri dosyalarının izlenebilirliğine de bağlıdır. Ay yerçekimi katsayı dosyası, topografya rasterı, albedo ürünü, SPICE kernel'i veya ST-LRPS eğitim bulutu değiştiğinde sonuçlar da değişebilir. Bu nedenle her koşunun çıktısı, yalnızca sayısal sonuçları değil, veri kaynağı bilgilerini de saklamalıdır.

Her görev koşusu için önerilen minimum izlenebilirlik paketi şudur:
\begin{itemize}
    \item görev başlangıç zamanı ve süre bilgisi,
    \item integratör yöntemi, toleranslar ve maksimum adım,
    \item kullanılan yerçekimi modeli yolu ve derece bilgisi,
    \item topografya/albedo/efemeris dosya yolları,
    \item etkin fizik modelleri,
    \item ST-LRPS kullanıldıysa run directory, config ve checkpoint bilgisi,
    \item Monte Carlo kullanıldıysa tohum ve dağılım parametreleri,
    \item çıktı üretim zamanı ve yazılım sürümü veya en azından çalıştırma klasörü.
\end{itemize}

Tekrar üretilebilirlik açısından en önemli ilke, ``aynı klasörü açan bir kişi aynı koşuyu anlayabilmelidir'' ilkesidir. Bu nedenle zaman damgalı çıktı klasörleri, kaydedilen JSON konfigürasyonları ve rapor üretim çıktıları yalnızca kullanım kolaylığı değil, mühendislik denetlenebilirliği sağlar.

\section{ST-LRPS İçin Kabul Kriterleri}
\label{sec:st_lrps_acceptance}

ST-LRPS modeli ana propagasyon hattına bağlanmadan önce yalnızca training loss değerine bakmak yeterli değildir. Eğitim loss'u düşmüş olsa bile model belirli irtifa bandında, belirli yönlerde veya OOD bölgelerde başarısız olabilir. Bu nedenle ST-LRPS için kabul kriterleri çoklu metrik üzerinden tanımlanmalıdır.

\subsection{Zorunlu Metadata Kontrolleri}

Bir ST-LRPS koşu klasörü kullanılmadan önce aşağıdaki koşullar sağlanmalıdır:
\begin{itemize}
    \item \texttt{config.json} ve \texttt{scaler.json} mevcut olmalıdır.
    \item \texttt{checkpoints/ckpt\_best.pt} veya \texttt{checkpoints/ckpt\_last.pt} bulunmalıdır.
    \item \texttt{central\_body} açıkça Ay olmalıdır.
    \item \(\mu\) ve referans yarıçapı Ay sabitleriyle uyumlu olmalıdır.
    \item \texttt{degree\_min}, \texttt{degree\_max} ve \texttt{target\_mode} eğitim ve değerlendirme verisiyle uyuşmalıdır.
    \item konum ölçekleme politikası Ay merkezini kaydırmamalıdır.
\end{itemize}

Bu koşullardan biri sağlanmıyorsa modelin sayısal olarak tahmin üretmesi yeterli değildir; model fiziksel olarak geçersiz kabul edilmelidir.

\subsection{Metrik Tabanlı Kabul}

ST-LRPS değerlendirmesinde minimum raporlanması gereken metrikler şunlardır:
\[
\mathrm{RMSE}_{a}
=
\sqrt{
\frac{1}{N}
\sum_{i=1}^{N}
\|\Delta\mathbf{a}_{\theta,i}-\Delta\mathbf{a}_{i}\|^2
},
\]
\[
\theta_i
=
\cos^{-1}
\left(
\frac{
\Delta\mathbf{a}_{\theta,i}\cdot\Delta\mathbf{a}_{i}
}{
\|\Delta\mathbf{a}_{\theta,i}\|\,\|\Delta\mathbf{a}_{i}\|
}
\right).
\]
Burada \(\mathrm{RMSE}_a\) ivme büyüklük hatasını, \(\theta_i\) ise vektörel yön hatasını gösterir. Buna ek olarak radyal ve cross-radial hata bileşenleri ayrı raporlanmalıdır:
\[
e_R = \mathbf{e}_a\cdot\hat{\mathbf{r}},
\qquad
e_C = \|\mathbf{e}_a-e_R\hat{\mathbf{r}}\|.
\]

Modelin kabulü için tek bir evrensel sayı verilmesi doğru değildir; çünkü görev irtifası, süre, gereken navigasyon hassasiyeti ve modelin kullanılacağı karar tipi farklıdır. Ancak pratik kabul yaklaşımı şu olmalıdır:
\begin{itemize}
    \item global RMSE düşük olmalı,
    \item irtifa kutularında hata tek bir bölgede patlamamalı,
    \item yön hatasının medyan ve p95 değerleri raporlanmalı,
    \item OOD alt ve üst bantlarda hata artışı anlaşılır olmalı,
    \item ST-LRPS ile klasik yüksek dereceli SH propagasyonu kısa doğrulama koşusunda karşılaştırılmalıdır.
\end{itemize}

\subsection{Training Loss Neden Tek Başına Yeterli Değildir?}

Training loss, modelin eğitim veri kümesine uyumunu gösterir. Ancak görev propagasyonu açısından kritik olan, modelin yeni konumlarda doğru ivme üretmesi ve bu hatanın zamanla yörüngeye nasıl yansıdığıdır. Bu yüzden validation ve test veri kümeleri eğitimden bağımsız olmalı; mümkünse farklı tohumlarla üretilmelidir.

Özellikle direction loss devreye geç girdiği için erken epoch'ta düşük görünen bir reference loss yanıltıcı olabilir. Bu nedenle checkpoint seçimi direction loss ramp tamamlandıktan sonra başlatılır. Böylece en iyi model seçimi, yalnızca potansiyel yüzeyi öğrenmiş ama yön davranışı olgunlaşmamış bir ağırlık setine takılmaz.

\section{Monte Carlo Güvenilirliği ve Senaryo Tasarımı}
\label{sec:mc_reliability}

Monte Carlo analizi, belirsizliklerin görev sonucu üzerindeki etkisini görmek için kullanılır. Ancak Monte Carlo'nun güvenilirliği, örnek sayısından önce dağılım varsayımlarına bağlıdır. Yanlış veya gerçekçi olmayan dağılımlar seçilirse 10.000 senaryo bile yanlış güven hissi verebilir.

Monte Carlo tarafında örneklenebilecek belirsizlikler üç gruba ayrılabilir:
\begin{itemize}
    \item \textbf{Başlangıç durumu belirsizliği:} konum, hız, orbital elemanlar ve epoch belirsizliği.
    \item \textbf{Uzay aracı parametre belirsizliği:} kütle, alan, \(C_D\), \(C_R\), area-to-mass oranı.
    \item \textbf{Model belirsizliği:} yerçekimi derece seçimi, ST-LRPS/classic backend seçimi, topografya sağlayıcısı, SRP/albedo/termal model bayrakları.
\end{itemize}

Monte Carlo sonuçları yorumlanırken yalnızca ortalama değer değil, dağılımın kuyruğu da incelenmelidir. Düşük olasılıklı fakat görev kaybına neden olabilecek çarpışma senaryoları, ortalama irtifa metriğinde görünmeyebilir. Bu nedenle raporda şu metrikler öncelikli olmalıdır:
\begin{itemize}
    \item impact probability,
    \item survival ratio,
    \item minimum altitude percentile değerleri,
    \item final altitude dağılımı,
    \item 95\% güven aralığı,
    \item en kötü senaryo örneklerinin parametreleri.
\end{itemize}

ST-LRPS Monte Carlo içinde kullanıldığında ek dikkat gerekir. Vekil model hız kazandırabilir; fakat modelin eğitim irtifa bandı dışına çıkan senaryolarda güvenilirliği azalabilir. Bu nedenle Monte Carlo sonrası yalnızca çarpışma oranı değil, senaryoların eğitim bandı içinde kalıp kalmadığı da analiz edilmelidir.

\section{Operasyonel Çalıştırma Prosedürü}
\label{sec:operational_workflow}

Projenin gerçek kullanımı için önerilen operasyonel akış aşağıdaki gibidir:

\begin{enumerate}
    \item \textbf{Veri kontrolü:} Yerçekimi, efemeris, topografya ve ST-LRPS klasörleri seçilir; preflight kontrolleri çalıştırılır.
    \item \textbf{Nominal koşu:} Önce kısa süreli ve düşük maliyetli bir görev koşusu yapılır.
    \item \textbf{Fizik modeli artırımı:} SRP, üçüncü cisim, albedo, termal, tides ve relativistik etkiler sırayla açılarak farklar incelenir.
    \item \textbf{Yüksek sadakat koşusu:} Nihai fizik bayraklarıyla uzun görev propagasyonu yapılır.
    \item \textbf{Rapor üretimi:} Zaman damgalı çıktı klasöründe grafikler ve PDF raporu oluşturulur.
    \item \textbf{Monte Carlo koşusu:} Belirsizlikler tanımlanır ve yeterli senaryo sayısı ile toplu analiz yapılır.
    \item \textbf{Sonuç inceleme:} Etki olasılığı, irtifa zarfı, final dağılım ve worst-case senaryolar incelenir.
    \item \textbf{Arşivleme:} Konfigürasyon, veri yolları, raporlar ve analiz çıktıları aynı run klasöründe korunur.
\end{enumerate}

Bu prosedür, kullanıcının doğrudan en karmaşık modele atlamasını engeller. Karmaşık model tek seferde çalıştırıldığında hata çıktısının kaynağını anlamak zorlaşır. Aşamalı çalıştırma, hem debugging hem de fiziksel yorum açısından daha güvenlidir.

\section{Sınırlılıklar ve Gelecek Geliştirmeler}
\label{sec:limitations_future_work}

Mevcut yapı yüksek sadakatli bir Ay simülasyon platformu sunar; ancak her model gibi sınırları vardır. Bu sınırların açıkça yazılması, çalışmanın zayıflığı değil güvenilirliğinin bir parçasıdır.

\subsection{Mevcut Sınırlılıklar}

\begin{itemize}
    \item ST-LRPS yalnızca eğitildiği irtifa, derece aralığı ve veri dağılımı içinde güvenilir kabul edilmelidir.
    \item GPU Monte Carlo yolu tüm fizik modellerini kapsamayabilir; karmaşık yüzey ve vekil model durumlarında CPU yoluna dönmek gerekebilir.
    \item Topografya verisi yoksa sabit yarıçap yaklaşımı kullanılır; bu, gerçek yüzey emniyeti ile aynı anlamda değildir.
    \item Efemeris kernel'leri ve zaman kapsamı dışına çıkıldığında interpolasyon güvenilirliği düşer.
    \item Radyal/cross-radial hata ayrımı, hız vektörü kullanılmadığında gerçek RTN ayrışımı değildir.
    \item Monte Carlo sonuçları, seçilen belirsizlik dağılımları kadar güvenilirdir.
\end{itemize}

\subsection{Gelecek Geliştirmeler}

Gelecekte sistemi daha güçlü hale getirmek için aşağıdaki geliştirmeler önerilir:
\begin{itemize}
    \item ST-LRPS için görev bazlı acceptance dashboard oluşturulması,
    \item Monte Carlo worst-case senaryolarının otomatik yeniden propagasyonu,
    \item topografya çarpışma olaylarında lokal eğim ve yüzey normali bilgisinin eklenmesi,
    \item ST-LRPS model belirsizliği için ensemble veya dropout tabanlı güven aralığı tahmini,
    \item GPU yolunun daha fazla fizik modelini kapsayacak şekilde genişletilmesi,
    \item veri dosyaları için hash tabanlı provenance kaydı,
    \item raporda otomatik hata bütçesi ve tolerans hassasiyeti tablosu üretilmesi.
\end{itemize}

\section{Operasyonel Kabul Checklist'i}
\label{sec:operational_checklist}

Bir sonuç görev kararı için kullanılmadan önce aşağıdaki kontrol listesi tamamlanmalıdır:

\begin{itemize}
    \item Başlangıç epoch değeri açık ve UTC/J2000 dönüşümü tutarlı mı?
    \item Yerçekimi modeli dosyası, derece ve referans yarıçapı beklenen değerlerle uyumlu mu?
    \item ST-LRPS kullanılıyorsa model klasörü Ay metadata'sı taşıyor mu?
    \item ST-LRPS evaluation raporunda test ve OOD metrikleri kabul edilebilir mi?
    \item Topografya çarpışma analizi gerçek LDEM verisiyle mi, sabit küreyle mi yapıldı?
    \item Integratör toleransı ve maksimum adım görev süresine uygun mu?
    \item Kuvvet bütçesi raporunda beklenmeyen sıfır veya aşırı büyük ivme bileşeni var mı?
    \item Monte Carlo örnek sayısı ve dağılım varsayımları raporda açık mı?
    \item Çıktı klasöründe config, log, grafik ve rapor dosyaları birlikte saklandı mı?
    \item Aynı koşu başka bir makinede veya başka bir oturumda yeniden üretilebilir mi?
\end{itemize}

Bu checklist, simülasyonun yalnızca ``tamamlandı'' durumuna değil, mühendislik açısından ``kabul edilebilir'' durumuna geçmesini sağlar.
